% IEEE conference draft. Original results and experimental workflows are unchanged.
\documentclass[conference]{IEEEtran}
\IEEEoverridecommandlockouts
\usepackage[T1]{fontenc}
\usepackage{cite}
\usepackage{amsmath,amssymb}
\usepackage{booktabs}
\usepackage{listings}
\usepackage{xcolor}
\usepackage[hidelinks]{hyperref}
\newcommand{\todo}[1]{\textcolor{red}{[TODO: #1]}}
\newcommand{\code}[1]{\nolinkurl{#1}}
\lstset{basicstyle=\ttfamily\scriptsize,breaklines=true,columns=fullflexible,keepspaces=true,frame=single}
\title{TRACE: Evidence-Grounded Issue Specifications for Regression Test Generation}
\author{
\IEEEauthorblockN{Mark Koehler}
\IEEEauthorblockA{DePaul University\\mkoehle6@depaul.edu}
\and
\IEEEauthorblockN{Hikaru McGuan}
\IEEEauthorblockA{DePaul University\\hmcguan@depaul.edu}
\and
\IEEEauthorblockN{Jason Torres}
\IEEEauthorblockA{DePaul University\\jtorr138@depaul.edu}
\thanks{Course project for SE 433: Software Testing and Quality Assurance, DePaul University.}
}
\begin{document}
\maketitle
\begin{abstract}
Issue-driven regression testing requires translating informal reports into executable behavioral checks. TRACE (Test Reproduction through Analysis, Checks, and Evidence) adds bounded extraction and validation of evidence-grounded issue specifications before a DeepSeek-Flash-powered coding agent generates native test patches. We compare this workflow with a shared native-test baseline on 26 assigned SWE-bench Verified instances using the official SWT-Bench unit-test evaluator. Baseline succeeds on 16 instances and TRACE on 17; all 26 patches from each workflow apply successfully. The paired outcomes are 15 shared successes, two TRACE-only successes, one baseline-only success, and eight shared unsuccessful outcomes. Mean generation time increases from 93.45 to 127.09 seconds. Six systematically selected case studies illustrate that issue-specification acceptance does not guarantee a test that correctly checks the reported behavior or compatibility with the fixed repository. Test-name attribution discrepancies, incomplete historical image records, diagnostic deviations from the protocol, incomplete per-method cost accounting, and single-run variability limit interpretation. The observed one-instance gain is preliminary evidence rather than a demonstrated general improvement.
\end{abstract}
\begin{IEEEkeywords}
regression testing, test generation, coding agents, SWT-Bench, issue specifications
\end{IEEEkeywords}
\section{Research Question and Motivation}
A regression test must encode the reported behavior, execute in the project's native suite, and distinguish buggy from corrected code. An agent can produce a plausible assertion that passes prematurely, fails for an unrelated reason, or becomes invalid after a legitimate API change. SWT-Bench evaluates repository-level generated tests against reference fixes and supplies complementary behavioral and coverage evidence~\cite{swtpaper,swtrepo}.

We use DeepSeek-Flash for test generation in both workflows. We ask: \emph{Does evidence-grounded issue-specification extraction and validation improve DeepSeek-Flash's issue reproduction, as measured by official SWT-Bench instance-level success, compared with direct report-to-test generation?} TRACE (Test Reproduction through Analysis, Checks, and Evidence) is motivated by the hypothesis that explicitly identifying a target, trigger, observed behavior, expected behavior, and suggested assertion will reduce unsupported assumptions. We measure reproduction success, not the number of unsupported assumptions, so the proposed mechanism is not directly tested. Validation is intended to filter the issue specification before generation; it does not independently prove test correctness. Our contribution is a bounded preprocessing workflow, preserved paired artifacts, and an empirical examination of when structured issue specifications help or fail.

\section{Method and Experimental Design}\label{sec:method}
\subsection{Two Workflows}
In this paper, ``DeepSeek'' means DeepSeek-Flash; recorded provider identifiers appear in Appendix~\ref{app:disclosure}.

\textbf{Baseline.} Our Python script, \code{controller.py}, reads the saved issue description from \code{report.txt}. It passes that text unchanged to mini-SWE-agent 2.4.6~\cite{mini} by calling \code{agent.run(task)}. The controller does not summarize or analyze the issue first. mini-SWE-agent combines the issue with the shared instructions for writing regression tests and sends these messages to DeepSeek. DeepSeek proposes shell commands to inspect source code, read existing tests, write a test, and run it. mini-SWE-agent executes those commands in a Docker container containing the buggy repository, then sends their output and exit status back to DeepSeek. This exchange repeats as the model examines results and adjusts its test. DeepSeek selects the actions; mini-SWE-agent carries them out and maintains the conversation. The controller saves the final submitted patch and the conversation record.

\textbf{TRACE.} Before starting that same agent process, the controller makes a separate request to DeepSeek to analyze only the original issue. It asks for five items: the target function, triggering input, observed behavior, expected behavior, and a suggested assertion. Each known item must include an exact supporting quote from the issue; an unknown item and its evidence must both be marked null. Python checks that the response has the required structure, that its quotes occur in the report, and that at least one item contains usable information. A second, separate DeepSeek request compares the extracted information with the original issue description. It checks whether the issue supports each item, whether any item contradicts the issue, and whether important information was left out. These checks assess the issue specification, not the generated test.

If either check rejects the issue specification, the controller sends the errors back with the original issue and requests another extraction. It permits up to three extraction attempts; each attempt that passes the Python checks is followed by a review, giving at most six preprocessing requests. When an issue specification is accepted, the controller appends it to the original issue and passes both to mini-SWE-agent. If all attempts are rejected, it passes only the original issue and records a fallback. Missing credentials and other operational errors stop the run rather than triggering this fallback. The agent then uses the same DeepSeek-driven inspection, test-writing, execution, and submission process as the baseline. The specification-extraction and review requests are separate from the agent's generation conversation; only the accepted issue specification is handed from one to the other.

The course guidelines require a same-model workflow comparison and official unit-test evaluation~\cite{guidelines}. The preserved run records support this procedure~\cite{artifacts}. Both methods use the same native-test instructions, buggy environment, and generation limits: 100 agent steps and a nominal \$0.50 agent budget, excluding TRACE preprocessing. A final model request can take spending above that nominal limit. The agent's submitted patch is exported unchanged. All 52 selected runs submitted patches, so no empty predictions occurred in the final evaluated set; the exporter retains non-submissions as empty predictions when they occur. Neither method receives the reference fix or evaluator-only results during generation or repair. Running a test on buggy code supplies feedback, but does not establish that it passes after the fix. An automatic classifier that distinguishes issue-related failures from syntax, import, or setup errors is future work, not an implemented controller step.

The original agent configuration asks for a standalone script that reproduces the bug. For this experiment, both methods instead receive the same instructions to add regression tests to the repository's existing test suite. Both also enable Bash's \code{pipefail} option, which preserves a failing test runner's exit status when its output is piped through another command. These shared changes make the outputs suitable for unit-test evaluation and keep execution failures visible. The original configuration file is preserved; the controller applies these changes in memory for each run. The comparison therefore measures the addition of TRACE's issue specification to this shared unit-test workflow.

\subsection{Assignment, Controls, and Evaluation}
The fixed assignment comprises 26 issues, 13 Easy and 13 Hard, across 12 Python repositories. These are historical assignment labels, not our measured difficulty classes. We use SWE-bench Verified revision \code{c104f840cc67f8b6eec6f759ebc8b2693d585d4a}. Both arms use the same saved reports and the same Docker image for each pair. For each issue, we compare one baseline run with one TRACE run. We reused three such pairs from the pilot phase; for the remaining issues, we ran the baseline first and TRACE second. Configuration and report hashes are checked. The three pilot pairs recorded Docker image names rather than unique image fingerprints, limiting our ability to verify their original image contents. We have one selected generation per arm per instance, with earlier failed attempts preserved. No explicit generation seed was set. We did not repeat each method on each issue, so we cannot determine how much the results would change across repeated runs. Each assigned instance is the unit of analysis, not each generated test method. We record problems such as missing API credentials, network errors, and environment setup failures separately from rejected issue specifications. Fixing these setup problems does not permit changing the experimental instructions. For each issue and method, we use the generated patch identified in the saved run-selection record. We do not compare evaluation scores from earlier attempts and then choose the best-performing patch.

After test generation, we evaluate the saved patches using the official SWT-Bench harness at commit \code{eee38624e36a1b39976fffa2bddbc4c3c754d293}. The evaluator uses unit-test mode with standard patch handling and coverage collection enabled. It runs up to two evaluation jobs at once, with a 30-minute timeout for each execution. Both methods use identical evaluation settings. Dataset reference fields and evaluation outputs remain in an evaluation-only directory. All 26 assigned instances remain in the success-rate and patch-applicability denominators for each method. After infrastructure failures, recovery reuses unchanged environment-build recipes or predictions and preserves unsuccessful attempts. Candidates are not selected based on evaluation outcomes.

\section{Results}
\subsection{Success, Applicability, and Paired Outcomes}
Table~\ref{tab:results} reports official outcomes over all assignments. Baseline succeeds on 16/26 instances (61.54\%); TRACE succeeds on 17/26 (65.38\%). Both methods have 26/26 applicable patches (100\%): applicability means that the submitted patch applies cleanly to the repository. Table~\ref{tab:all} in Appendix~\ref{app:outcomes} lists every paired outcome and coverage value. Appendix~\ref{app:repro} provides reproducibility details, and Appendix~\ref{app:disclosure} records team contributions and AI-tool use. All 52 evaluation reports were produced---26 per method. No issues were removed from the totals because of technical problems. The evaluator reported no instance-level errors. Initial image timeouts and setup-check failures were resolved before final evaluation.

An evaluator exit code of zero indicates completed evaluation, not that every reproduction succeeded. The harness compares runs with the generated patch against a base comparison without that patch. It requires (i) additional fail-to-pass tests, (ii) preservation of the base comparison's fail-to-pass tests, and (iii) no additional tests that fail in both versions or change from passing to failing. Therefore, a patch can introduce a useful fail-to-pass test yet be unsuccessful because another assertion introduces a lasting regression. This stricter instance-level criterion explains some disagreements better than counting isolated failing assertions.

\begin{table}[t]
\caption{Official Evaluation and Selected Generation Records}
\label{tab:results}
\centering\small
\begin{tabular}{@{}p{0.48\columnwidth}rr@{}}
\toprule
Measure & Baseline & TRACE\\
\midrule
Official successes / assigned & 16/26 & 17/26\\
Success rate & 61.54\% & 65.38\%\\
Applicable / assigned & 26/26 & 26/26\\
Applicability rate & 100\% & 100\%\\
Mean coverage delta (cases with coverage) & 0.567364 & 0.602731\\
Coverage availability & 25/26 & 25/26\\
Mean coverage delta, successful cases & 0.691899 & 0.726739\\
Successful cases with coverage & 16 & 17\\
Mean generation time (seconds) & 93.45 & 127.09\\
Agent-only cost (USD) & 0.379072 & 0.383181\\
Evaluation time (seconds) & 3425.92 & 1549.86\\
\bottomrule
\end{tabular}
\end{table}

The paired outcomes are 15 instances where both succeed, one where only the baseline succeeds, two where only TRACE succeeds, and eight where neither succeeds (Table~\ref{tab:paired}). TRACE succeeds where the baseline does not on pytest 7432 and Django 15098; the baseline alone succeeds on Django 13344. The net difference is one instance, or 3.85 percentage points. The methods produced different success outcomes on only three issues: two favored TRACE and one favored the baseline. An exact two-sided McNemar test on the two TRACE-only and one baseline-only successes gives $p=1.0$. This small difference does not provide sufficient statistical evidence that either method performs better, nor does it establish that they perform equally.

\begin{table}[t]
\caption{Paired Official Reproduction Outcomes}
\label{tab:paired}
\centering\small
\begin{tabular}{lrr}
\toprule
 & TRACE succeeds & TRACE unsuccessful\\
\midrule
Baseline succeeds & 15 & 1\\
Baseline unsuccessful & 2 & 8\\
\bottomrule
\end{tabular}
\end{table}

TRACE accepts 22 issue specifications and falls back after three attempts on Django 15098, Django 16631, Sphinx 8551, and Django 13810. These fallbacks remain TRACE outcomes. Django 15098 is a TRACE-only success despite fallback, so that gain cannot be credited directly to the accepted issue specification. Django 11163 and Flask 5014 have byte-identical baseline/TRACE patches. TRACE accepts the issue specification for Django 13344 on its second extraction attempt, without fallback. Overall counts therefore describe the complete workflow, including fallback and sampling variation, rather than establishing whether specification extraction caused the difference.

\subsection{Coverage, Cost, and Runtime}
We report the official field \code{coverage_delta_pred}, the fraction of executable lines changed by the reference fix whose execution counts increase relative to the harness base comparison. This fraction is unitless. It is not whole-repository coverage. Arithmetic means exclude unavailable values and are not weighted by repository size or changed-line count. Overall means include all cases with available coverage, regardless of success. Success-only means use each method's successful cases with available coverage, so the two methods are averaged over different subsets.

Sphinx 8721 has N/A coverage in both arms because the harness finds no executable changed lines in its coverage inputs. N/A is excluded, not replaced with zero. The per-instance coverage values and aggregation counts are indexed in \code{evidence.json}; the metric implementation is in \code{src/grading.py} at the recorded SWT-Bench harness commit. A zero coverage delta means that no counted changed line has a positive execution-count increase relative to the base comparison; it does not mean that those lines were never executed. Success uses test outcomes rather than this coverage measure, so TRACE can succeed on Django 15098 and SymPy 24539 with zero delta. Baseline Django 15098 also has zero delta, but is unsuccessful. Identical Django 11163 patches receive different coverage deltas, showing that these measurements can vary even when submitted content matches.

Agent-reported costs for the selected runs total \$0.379072 for baseline and \$0.383181 for TRACE. TRACE additionally makes 69 preprocessing requests, recording 49,789 input and 173,755 output tokens. These costs exclude preprocessing and are not the total experiment cost. We report the recorded preprocessing token usage separately without assigning it a dollar amount. 

DeepSeek billing exports for October 18--November 16, 2026 show total charges of \$0.959976 (rounded) across the exported period. These charges are not limited to the selected 52 generation runs and cannot be separated reliably into baseline and TRACE totals from the daily billing records alone. They are reported separately, not added to the agent-reported costs, because adding them would double-count charges for the selected runs. 

Mean generation time recorded by the controller increases by 33.64 seconds, approximately 36\%, including TRACE preprocessing but excluding image preparation. Evaluation takes 57.10 minutes for baseline and 25.83 minutes for TRACE. TRACE runs second and reuses comparison caches, so evaluation timing is not evidence of workflow efficiency.

\section{Test Quality Analysis}\label{sec:test-quality}
Both the baseline and TRACE workflows instruct the agent to include meaningful assertions and inspect failures for syntax, import, setup, or issue-related causes. Execution records support manual review, but we did not implement an independent automatic checker for assertion quality or failure relevance.

We choose six distinct cases systematically: sort full instance IDs within each category, take the first two shared successes, first two shared failures, then first two disagreements, without duplication. This yields Astropy 14508, Django 11163, Django 12155, Django 13513, Django 13344, and Django 15098. Each category contains at least two cases, so no additional cases are needed. Appendix~\ref{app:cases} discusses the issue, generated check, both execution states, semantic relevance, and method implication for each.

The successful cases check clear expected behaviors: writing decimal numbers compactly in astronomy file headers without truncating comments, and returning an empty dictionary when no fields are requested. Both methods capture these behaviors; additional structure offers no demonstrated advantage there. The failures expose different weaknesses. For Django 12155, both methods produced assertions that exposed the indentation problem before the fix. However, the reference fix removed a helper function that the existing test module imported. This prevented the module from loading afterward, so neither method demonstrated a passing test on the fixed code. For Django 13513, both methods added tests checking that Django's debug report hides an earlier exception explicitly suppressed with \code{raise ... from None}. These tests passed before and after the fix because that behavior already worked. They did not exercise the scenario addressed by the official fix---an exception without a traceback---and therefore did not reproduce the bug. An accepted issue specification does not guarantee compatibility with the fixed code or the ability to distinguish buggy from fixed code.

The disagreement cases illustrate that an added persistent failure can make an instance unsuccessful even when new tests go from failing to passing. They also show that parsed test labels can differ from raw execution logs. Django 13344 shows that validating an issue specification does not guarantee a correct test setup: TRACE reproduced a similar symptom in a different middleware arrangement, but its test still failed after the fix. For Django 15098, both methods generated URL tests that failed before the fix and passed afterward. However, the baseline also added a language-lookup assertion that continued failing after the fix, causing its overall evaluation to fail. TRACE avoided that additional persistent failure and succeeded. Because TRACE used the original report after its issue specification failed validation, this result does not demonstrate a benefit from the extracted issue specification. Several official transition labels do not align cleanly with raw test-name lines; we report the official score and raw observations separately rather than silently correcting either.

\section{Threats to Validity}
Several factors limit our ability to attribute the observed improvement to TRACE. We selected one run per method for each issue, reused three pilot pairs, and ran the baseline before TRACE for the remaining issues. We did not measure variation across repeated runs or changes in the hosted model. The same model also extracted and reviewed the issue specification, so both steps could share the same misunderstanding.

Some recorded agent actions require additional review. Six selected runs were flagged for inspecting repository history: baseline on Django 13513, Django 15098, and Django 16631; TRACE on Seaborn 3187, Sphinx 8548, and Sphinx 8638. These are command-review flags, not proof that a reference fix was accessed or used. Other diagnostic actions included a temporary plugin that changed software behavior and stand-in code for missing dependencies. For Sphinx 8621, the baseline agent created a temporary pytest plugin that replaced the keyboard-formatting function in memory to try out its own proposed fix; passing tests under that plugin were not results from the official reference fix. For Sphinx 8638, the TRACE agent supplied a temporary replacement for the missing \code{roman} dependency, whose \code{toRoman()} function returned ordinary number strings instead of Roman numerals. This allowed diagnostic tests to run, but did not reproduce the real dependency's behavior. The submitted patches contain only test changes, but the patches alone do not establish that all actions during generation followed the experiment's restrictions.

The reported results depend on which tests the evaluator runs and how it reads their output. Some test names could not be matched between the before-fix and after-fix logs, and some recorded before/after transitions belonged to existing tests rather than newly generated ones. These issues can affect the reported success counts. Coverage measures execution of lines changed by the reference fix, not whether an assertion correctly represents the issue. We report the official results for all 26 issues and document discrepancies found in the raw logs. For three pilot pairs, we recorded Docker image names but not unique image fingerprints, so we cannot fully verify their original contents.

Our findings apply to a small set of assigned Python issues. Neither the Easy/Hard labels nor TRACE's one-instance advantage indicates how either method would perform on other projects. Provider billing totals are available for the exported period, but the records do not separate charges by workflow or selected run, so they cannot establish which method is more cost-effective.

A follow-up study should specify the number of repeated runs and how rule violations will be handled before reviewing any results. Human reviewers could also independently check whether each test captures the reported behavior and whether its failure comes from the bug rather than an environment problem. Measuring how often reviewers agree would provide an additional check beyond DeepSeek reviewing its own issue specification.

\section{Conclusion}
TRACE achieves 17/26 official successes versus 16/26 for the baseline, with 100\% patch applicability and increased mean generation time. The paired evidence supports a modest observed difference, not a reliable advantage for TRACE. Future work should repeat generation under recorded seeds where supported, compare extraction with and without review, check how raw test results are converted into parsed transitions, and add a shared quality checker for issue relevance and setup failures. Such a checker should initially assess both methods without using evaluator-only results to repair tests. Any later analysis of how the findings change under alternative assumptions should supplement the full assigned-set results rather than discard unfavorable cases.

\begin{thebibliography}{00}
\bibitem{swtpaper} N. M\"undler, M. N. M\"uller, J. He, and M. Vechev, ``SWT-Bench: Testing and validating real-world bug-fixes with code agents,'' in \emph{Advances in Neural Information Processing Systems}, 2024, doi: 10.48550/arXiv.2406.12952.
\bibitem{swtrepo} SWT-Bench contributors, ``SWT-Bench evaluation harness,'' source revision \code{eee38624e36a1b39976fffa2bddbc4c3c754d293}. [Online]. Available: \url{https://github.com/logic-star-ai/SWT-Bench}
\bibitem{mini} mini-SWE-agent contributors, ``mini-SWE-agent,'' version 2.4.6. [Online]. Available: \url{https://github.com/SWE-agent/mini-swe-agent}
\bibitem{guidelines} DePaul SE Team, ``Final project: Improving DeepFlash test generation with mini-SWE-agent and SWT-Bench,'' secs. 6, 8, 11, and 12, accessed Nov. 13, 2026. [Online]. Available: \url{https://github.com/DePaul-SE-Team/SE443-333-final-project-guideline}
\bibitem{artifacts} TRACE project, ``Generation and official evaluation artifacts,' Nov. 2026. [Online]. Available: \url{https://github.com/mkoehle6/trace-swt-bench}. \todo{Add updated commit-pinned evaluation artifact snapshot link.}
\end{thebibliography}

\appendices
\section{Six Distinct Case Studies}\label{app:cases}
The cases below use the selection rule in Section~\ref{sec:test-quality}. ``Success'' denotes the official instance-level result. Raw before-fix and after-fix observations are reported separately when the log parser assigns results to different test names. Source paths and complete transition sets are indexed in \code{evidence.json} beside the paper.

We assess TRACE's generated tests using four properties: \textbf{P1, meaningful assertions}: the test explicitly checks an expected outcome; \textbf{P2, issue-related failure}: the test fails on buggy code because of the reported problem, not a typo, import error, or setup failure; \textbf{P3, passes after the fix}: the same test executes and passes after the reference fix; and \textbf{P4, requirement alignment}: the test uses a setup and assertions that represent the behavior requested in the issue. These are retrospective assessments from submitted patches and saved execution evidence, not outputs of an implemented automatic quality checker. P1--P4 are scored only for TRACE to assess the proposed workflow; baseline behavior is described for comparison but was not separately scored with this rubric. A passing post-fix test alone does not establish issue reproduction. ``Partial'' indicates that only part of the reported scenario is represented; ``not established'' indicates insufficient evidence, rather than a confirmed negative result.

\subsection{Shared Success: Astropy 14508}
\textbf{Issue.} FITS (Flexible Image Transport System) is a file format used in astronomy. Its header stores metadata in 80-character records called cards, containing a keyword, a value, and an optional comment. The HIERARCH convention allows keyword names longer than the usual eight characters, while retaining the 80-character record limit.\footnote{Astropy FITS header documentation: \url{https://docs.astropy.org/en/stable/io/fits/usage/headers.html}.}

In this issue, a decimal value is written with unnecessarily many digits, leaving too little space for the comment and causing it to be cut off. \textbf{Generated checks.} Both methods construct the reported card using the value 0.009125 and compare its serialized representation with the compact expected string. Baseline additionally checks the stored value and comment. \textbf{Before.} Both new tests fail with a VerifyWarning about comment truncation. \textbf{After.} Each passes, and each official report records one added fail-to-pass test without added persistent or pass-to-fail failures. \textbf{Relevance.} The tests check that the number is written compactly and the full comment is preserved, directly matching the problem described in the bug report. \textbf{Implication.} Both methods generate a native test that checks the expected behavior; TRACE's accepted issue specification does not produce an observable success advantage.
\begin{lstlisting}[language=Python,caption={Shared Astropy oracle, condensed from the submitted tests.}]
c = fits.Card("HIERARCH ESO IFM CL RADIUS", 0.009125,
              "[m] radius arround actuator to avoid")
assert str(c) == _pad(
    "HIERARCH ESO IFM CL RADIUS = 0.009125 / "
    "[m] radius arround actuator to avoid")
\end{lstlisting}
The comment spelling \code{arround} is reproduced verbatim from the issue example.

\textbf{TRACE quality assessment.} P1 met: the test explicitly compares the serialized card with the compact expected representation. P2 met: the buggy run encounters the reported comment-truncation warning while constructing the card, before reaching the final comparison; this is an issue-related failure, not a setup error. P3 met: the test passes after the fix. P4 met: the chosen value and comment directly exercise the reported formatting problem.

\subsection{Shared Success: Django 11163}
\textbf{Issue.} An explicitly empty fields list should produce an empty dictionary. \textbf{Generated checks.} Both workflows add the same test method to an existing test class that uses the database.
\begin{lstlisting}[language=Python,caption={Identical submitted Django 11163 addition.}]
def test_empty_fields(self):
    item = ColourfulItem.objects.create()
    self.assertEqual(model_to_dict(item, fields=[]), {})
\end{lstlisting}
\textbf{Before.} On buggy code, \code{model_to_dict(item, fields=[])} returns a dictionary containing model fields instead of an empty dictionary, so the assertion fails. \textbf{After.} The added test passes after the reference fix, and the official evaluator marks both workflows successful. It lists two fail-to-pass test entries, but only one is the newly generated test; the other is an existing many-to-many test. The reported count therefore does not mean that either agent wrote two tests. \textbf{Relevance.} The assertion checks exactly the behavior requested in the issue and uses a model already available in Django's test suite. \textbf{Implication.} Both workflows produced the same test patch, so TRACE's additional issue-specification step showed no benefit for this clear report in the observed runs. Nevertheless, the harness reported coverage deltas of 0.5 for baseline and 1.0 for TRACE. Because the patches are identical, patch differences cannot explain these coverage results.

\textbf{TRACE quality assessment.} P1 met: \code{assertEqual} checks for an empty dictionary. P2 met: the assertion fails because buggy code returns model fields for an explicitly empty field list. P3 met: the added test passes after the fix, as shown in the raw execution evidence. P4 met: the input and expected result directly match the issue. Parser attribution to an additional existing test does not change this assessment of the new test.

\subsection{Shared Failure: Django 12155}
\textbf{Issue.} When a docstring starts with text rather than a blank line, Django leaves extra indentation on later lines. This can cause errors when the documentation is displayed. \textbf{Generated checks.} Both workflows add assertions checking that indentation is removed correctly and that the docstring is split into the expected parts. TRACE also checks that the displayed documentation does not contain a formatting-error message from docutils, the documentation-processing library. \textbf{Before.} The saved execution logs show two new assertion failures for each workflow. These failures expose incorrect indentation handling; TRACE also detects a documentation-formatting error. \textbf{After.} The reference fix removes the \code{trim_docstring} helper, but the existing test module still imports it. That import fails before the runner can execute the newly added tests. The evaluator therefore cannot match those tests across the before-fix and after-fix runs, records no added fail-to-pass tests, and marks both workflows unsuccessful. \textbf{Relevance.} The assertions target the reported problem, but the tests rely on a helper that is no longer available after the fix. \textbf{Implication.} Exposing the bug before the fix is only part of a successful reproduction: the same test must also run and pass on the fixed code. Neither workflow demonstrated that here. Validating the issue specification did not prevent this test-compatibility problem, and the reference fix was not available to the agents during generation.

\textbf{TRACE quality assessment.} P1 met: the generated tests check whitespace handling and reject rendering-error output. P2 met for the observed new failures: retained indentation and rendering diagnostics concern the reported problem. P3 not met in this evaluation: the removed helper prevents the module from loading, so no passing post-fix execution of the new tests is demonstrated. P4 met for the intended indentation and rendering requirements, but reliance on the helper makes the submitted tests incompatible with the fixed repository.

\subsection{Shared Failure: Django 13513}
\textbf{Issue.} The report states that Django's debugging error page displays an earlier exception even when \code{raise ... from None} explicitly requests that it be hidden. This concerns Python's exception-context suppression. \textbf{Generated checks.} Both workflows create one exception while handling another and use \code{raise ... from None} to request that the earlier exception not be displayed. They then check the exception information collected for the debug report. \textbf{Before and after.} The new tests pass on both buggy and fixed code because Django already handles the particular scenario they check. Neither workflow produces a new test that fails before the fix and passes afterward, so both receive an unsuccessful result. A separate failure in an existing debug-template test is not a failure of these newly generated checks. \textbf{Relevance.} The tests address exception hiding, which is related to the report, but miss the specific case addressed by the official fix: an exception without a traceback, meaning that it carries no record of where it was raised. This detail was identified in our post-hoc review of the evaluation results and was not supplied to the agents for generation or repair. \textbf{Implication.} A test can contain sensible assertions about the right general topic yet fail to exercise the actual bug. TRACE's accepted issue specification did not prevent that outcome. Both methods receive a coverage delta of 0.647059 for this issue despite failing to reproduce it; executing relevant code is not the same as checking the broken behavior. The baseline agent also inspected repository history, which requires review against the experiment's restrictions.

\textbf{TRACE quality assessment.} P1 met: explicit assertions check the exception information. P2 not met: the generated tests pass on buggy code instead of exposing the bug. P3 met as an execution observation: the tests pass after the fix; because they also pass before it, this does not show successful reproduction. P4 partial: the tests check related exception-hiding behavior but omit the exception-without-a-traceback scenario. They therefore cannot distinguish the buggy version from the fixed version.

\subsection{Disagreement: Django 13344}
\textbf{Issue.} Middleware is code that processes a web request or response around a Django view. In the reported bug, the first middleware's \code{process_response()} method receives a coroutine, an asynchronous operation that still needs to be awaited, instead of a completed \code{HttpResponse}.

\textbf{Generated checks.} Both workflows add a helper middleware that records the object passed to \code{process_response()}. They then assert that the recorded object is an \code{HttpResponse}. The baseline places its helper before Django's existing \code{SecurityMiddleware}, following the report's example. TRACE instead places its helper before a new async-only middleware written by the agent. TRACE chose that setup after temporary experiments showed that it produced the reported coroutine symptom. Its issue specification had been accepted on the second extraction attempt, so this was not a fallback run.
\begin{lstlisting}[language=Python,caption={Expected response type checked by both workflows, simplified.}]
self.assertIsInstance(response, HttpResponse)
\end{lstlisting}

\textbf{Before.} Both new tests fail because their recording helpers receive a coroutine rather than an HttpResponse. \textbf{After.} The baseline test passes after the reference fix. TRACE's test still receives a coroutine and fails. The official evaluator therefore marks the baseline successful and TRACE unsuccessful.

\textbf{Reading the evaluation records.} The baseline's raw logs show the new test failing before the fix and passing afterward. However, the evaluator's log parser labels an existing test, \code{test_get_asgi_application}, as the fail-to-pass entry and cannot match the new test name between runs. We report the official outcome while noting that its test-name attribution differs from the raw logs.

\textbf{Relevance.} Both tests check the expected response type, but they use different middleware setups. TRACE reproduced a similar symptom without demonstrating that its setup exercised the same underlying bug. Its continued failure after the fix does not by itself identify the precise cause. \textbf{Implication.} Checking that the issue specification is supported by the report does not guarantee that the agent will build a faithful executable test. The response assertion and the surrounding setup both matter.

\textbf{TRACE quality assessment.} P1 met: the test explicitly checks for an HttpResponse. P2 not established: it fails with the reported symptom, but the evidence does not establish that the same underlying bug caused it. P3 not met: the test still fails after the reference fix. P4 partial: the expected response type matches the report, but the custom middleware setup differs from its example and is not shown to reproduce the same scenario.

\subsection{Disagreement: Django 15098}
\textbf{Issue.} Django does not correctly recognize language codes that specify both a writing system and a region. For example, \code{en-latn-us} means English, written using the Latin alphabet, for the United States. A valid web address using this language prefix can return a 404 Not Found response instead of the requested page. \textbf{Generated checks.} Both workflows add tests checking that the language-prefixed address works and that Django recognizes the language code. The baseline adds a language-lookup assertion to an existing test. TRACE creates a separate lookup test that also checks codes with mixed uppercase and lowercase letters. Before generation, TRACE tries three times to extract an issue specification that passes validation. None is accepted, so the agent receives only the original report, just as in the baseline. This run still counts as a TRACE result. Its recorded generation time includes the unsuccessful preprocessing attempts, and their API charges contribute to workflow spending; the agent-only cost figures exclude those requests.
\textbf{Before.} The URL tests and language-lookup assertions generated by both workflows fail. \textbf{After.} Both URL checks pass. The existing lookup test modified by baseline still fails, while TRACE's new lookup test passes in the raw log. Both suites retain additional existing failures or errors. The official baseline report records two added fail-to-pass entries but also an additional persistent failure relative to base, yielding an unsuccessful result. TRACE records three added fail-to-pass entries without an added persistent failure and succeeds. The parser cannot match TRACE's new lookup test name between runs, even though the raw after-fix log records it as passing.
\textbf{Relevance.} The URL tests directly check the reported behavior, while the additional language-lookup assertions check different aspects of language-code handling. \textbf{Implication.} The harness checks whether the generated patch introduces additional failures relative to its base comparison, not simply whether a new test changes from failing to passing. This TRACE-only success occurred without an accepted issue specification and therefore cannot demonstrate that an accepted issue specification caused the gain. Baseline's history inspection remains a separate limitation.

\textbf{TRACE quality assessment.} P1 met: explicit checks cover language-prefixed URLs and language lookup. P2 met for the observed generated checks: buggy execution fails on the language-code behavior rather than a syntax or setup error. P3 met in the inspected raw logs: the URL checks and new lookup method pass after the fix; parsed attribution of the lookup method remains uncertain. P4 met for the reported language/script/region behavior. TRACE receives official success, but because preprocessing fell back to the original report, these properties cannot be credited to an accepted issue specification.

\onecolumn
\section{Full Paired Outcomes}\label{app:outcomes}
Table~\ref{tab:all} retains the fixed assignment order used in the saved instance list, rather than sorting by repository or result. Success and Unsuccessful denote official instance-level outcomes.
\begin{table*}[t]
\caption{All Assigned Instances: Official Outcomes and Coverage Delta}
\label{tab:all}
\centering\scriptsize
\begin{tabular}{lccrr}
\toprule
Instance & Baseline & TRACE & Delta baseline & Delta TRACE\\
\midrule
\code{pydata__xarray-3677} & Success & Success & 1.000000 & 1.000000\\
\code{django__django-14089} & Success & Success & 0.500000 & 0.500000\\
\code{django__django-11163} & Success & Success & 0.500000 & 1.000000\\
\code{sympy__sympy-24539} & Success & Success & 0.285714 & 0.000000\\
\code{scikit-learn__scikit-learn-14894} & Success & Success & 0.416667 & 0.416667\\
\code{sphinx-doc__sphinx-9711} & Success & Success & 0.583333 & 0.583333\\
\code{astropy__astropy-14508} & Success & Success & 0.785714 & 0.785714\\
\code{pytest-dev__pytest-7432} & Unresolved & Success & 1.000000 & 1.000000\\
\code{pallets__flask-5014} & Success & Success & 1.000000 & 1.000000\\
\code{matplotlib__matplotlib-13989} & Success & Success & 0.500000 & 1.000000\\
\code{psf__requests-2931} & Success & Success & 0.750000 & 0.750000\\
\code{pylint-dev__pylint-6903} & Success & Success & 1.000000 & 1.000000\\
\code{mwaskom__seaborn-3187} & Success & Success & 1.000000 & 0.666667\\
\code{django__django-13513} & Unresolved & Unresolved & 0.647059 & 0.647059\\
\code{sphinx-doc__sphinx-8548} & Unresolved & Unresolved & 0.000000 & 0.000000\\
\code{django__django-13344} & Success & Unresolved & 0.096774 & 0.000000\\
\code{django__django-15098} & Unresolved & Success & 0.000000 & 0.000000\\
\code{django__django-12155} & Unresolved & Unresolved & 0.466667 & 0.466667\\
\code{sphinx-doc__sphinx-8638} & Unresolved & Unresolved & 1.000000 & 1.000000\\
\code{django__django-15252} & Unresolved & Unresolved & 0.000000 & 0.600000\\
\code{django__django-16631} & Success & Success & 0.652174 & 0.652174\\
\code{sphinx-doc__sphinx-8551} & Unresolved & Unresolved & 0.000000 & 0.000000\\
\code{django__django-13810} & Success & Success & 1.000000 & 1.000000\\
\code{sphinx-doc__sphinx-8721} & Unresolved & Unresolved & N/A & N/A\\
\code{sphinx-doc__sphinx-8621} & Unresolved & Unresolved & 0.000000 & 0.000000\\
\code{matplotlib__matplotlib-25311} & Success & Success & 1.000000 & 1.000000\\
\bottomrule
\end{tabular}
\end{table*}


\twocolumn
\section{Reproducibility and Outstanding Information}\label{app:repro}
The host computer is an ASUS Zenbook UX8402VV with a 13th-generation Intel Core i9-13900H processor, 20 logical processors, and approximately 32 GB of RAM. These specifications were verified after the experiment; the saved run records do not independently establish the historical hardware configuration. The workflow used Ubuntu/WSL and Docker, but the CPU and memory allocations available to those environments during execution were not recorded. DeepSeek inference used the hosted API rather than running locally on this hardware.

The evaluation directory is \code{evaluation_results/trace-assigned-26-unit-20261113-122920/}; final selected generation records are in \code{runs/assigned-pairs-20261113-002517/completion-pass-v6/}. Predictions, prompts, patches, trajectories, manifests, image references, workflow hashes, and official pre/post logs are preserved. The dataset revision and harness commit appear in Section~\ref{sec:method}; \code{evidence.json} indexes exact local evidence. Evaluation inputs contain reference fields and must remain outside generation context. No report preparation step modifies either workflow or prediction.

The nominal agent-cost limit excludes preprocessing and does not guarantee a strict cap on spending. Selected preprocessing metrics record 69 requests, 49,789 input tokens, and 173,755 output tokens, including reasoning output reported by the provider. The supplied provider exports, \code{amount-2026-10-18_2026-11-16.csv} and \code{cost-2026-10-18_2026-11-16.csv}, record 1,211 requests for \code{deepseek-flash}, 1,118,057 cache-miss input tokens, 20,423,295 cache-hit input tokens, and 644,720 output tokens. Total billed cost is \$0.959975793. Summing each token quantity multiplied by its listed unit price reproduces both the daily charges and the total exactly. Nonzero charges occur on November 3 and November 10--13. This export-period total may include pilot runs, retries, and other requests; its rows do not identify the method, instance, or individual generation run. We therefore do not allocate it between baseline and TRACE or treat it as the cost of the selected runs alone. No explicit generation seed was configured. The submission date is November 17, 2026. The submission artifacts are archived at commit \code{a5344b0abf54dcb9674f70261ad1a66d938f8a3f}: \url{https://github.com/mkoehle6/trace-swt-bench/tree/a5344b0abf54dcb9674f70261ad1a66d938f8a3f}. The evaluation results are preserved in its \code{evaluation_results/} directory.

\section{Contributions and AI Disclosure}\label{app:disclosure}
Mark Koehler, Hikaru McGuan, and Jason Torres completed this project as a collaborative team effort, sharing responsibility for the research, implementation, evaluation, analysis, and presentation of the work. Author names and email addresses appear on the title page.

The following tools were used in the experiment or supporting project work. Recorded model names are reported as identifiers, not as proof of an unchanged provider backend.

\begin{itemize}
\item \textbf{mini-SWE-agent 2.4.6:} the agent framework that manages the generation conversation, executes shell commands in Docker, and collects submitted test patches. It is not itself a language model. The dependency is pinned in \code{requirements.txt}, and the controller checks this version before launch.
\item \textbf{DeepSeek-Flash:} generated candidate tests and supplied issue extraction and semantic-review responses. The generation configuration records \code{deepseek/deepseek-v4-flash}; the controller's default direct-API analysis identifier is \code{deepseek-flash}. Saved provider responses include the returned model identifier and, where supplied, a system fingerprint. These records do not establish an exact immutable model release for every request.
\item \textbf{OpenAI Codex:} assisted with implementation and workflow orchestration, artifact inspection, evaluation support, analysis, workflow diagrams for presentation slides, and manuscript drafting and revision. The historical Codex application/CLI version and underlying model identifier were not consistently recorded and cannot be specified reliably from the experiment artifacts. A current installation version would not establish the version used for earlier work.
\item \textbf{ChatGPT (GPT-6.1):} supported initial project planning, workflow discussions, troubleshooting, and explanations. The GPT-6.1 model designation was confirmed by the project owner. ChatGPT was not an additional benchmark test-generation arm.
\item \textbf{Claude Sonnet 5.5:} assisted with reviewing the repository, drafting presentation slides, and proofreading the paper; the team reviewed and revised the resulting material. The model designation was confirmed by the project owner.
\end{itemize}

The team confirms that all AI tools used for this project have been disclosed above. This paper does not claim that humans independently validated every generated assertion or verified every action against the experiment's rules. AI-assisted project work and writing are disclosed separately from the evaluated test-generation methods.

\end{document}
